\documentclass[12pt]{article}
\usepackage[utf8]{inputenc}
\usepackage[T1]{fontenc}
\usepackage{hyperref}
\title{Tugen Language Resources for K-12 Curriculum Development}
\author{Compiled by Claude Sonnet 5 (claude-sonnet-5) for Andr\'as Kornai}
\date{}
\begin{document}
\maketitle

\section{Introduction}
Tugen (ISO 639-3 tuy; Glottocode tuge1241) is a Kalenjin (Nilotic) language spoken by about 140,000 people in the Rift Valley of western Kenya. Dialects reported include Lembus, Arror and Samor. Search results turned up no Tugen-specific school materials.

\section{Primary Educational Resources}
None were found specifically for Tugen.

\section{Linguistic and Reference Materials}
\begin{itemize}
\item Wikipedia entry on Tugen: \url{https://en.wikipedia.org/wiki/Tugen_language}
\item Wikipedia entry on the Kalenjin languages: \url{https://en.wikipedia.org/wiki/Kalenjin_languages}
\item Omniglot page on Tugen writing, which attributes an alphabet chart to Wolfram Siegel. I did not retrieve its contents: \url{https://omniglot.com/writing/tugen.htm}
\item Jerono, a University of Edinburgh Nilotic blog PDF. The search summary describes a study of Kalenjin orthography, with vowel and tone distinctions not captured by the Kiswahili-based spelling. I did not verify that it concerns Tugen specifically: \url{https://blogs.ed.ac.uk/nilotic/wp-content/uploads/sites/3585/2021/05/Jerono.pdf}
\item Creider and Creider, \emph{A Dictionary of the Nandi Language}. This is a dictionary of Nandi, a related Kalenjin language, not Tugen: \url{https://www.koeppe.de/titel_a-dictionary-of-the-nandi-language}
\end{itemize}

\section{Religious and Community Resources}
None were found.

\section{Digital and Audio Resources}
Rosetta Project (Long Now) entry for Tugen: \url{https://rosettapanglossia.longnow.org/index.php?title=Tuy}

\section{Curriculum Development Context}
No Tugen dictionary, primer or school programme was found. The orthography is reported as unsettled, since the existing Kiswahili-based Kalenjin spelling does not mark vowel distinctions or tone. Curriculum work would need to begin with orthography and lexicon development, possibly building on Nandi and Kipsigis materials. This dossier is short because little was found.
\end{document}